\documentclass[10pt,a4paper]{amsart}
\usepackage{amsmath,amssymb,amsthm,graphicx}
% For commutative diagrams we use Xy-pic
\usepackage[all]{xy}

% We use 2cell for 2-commutative diagrams.
\xyoption{2cell}
% Theorem environments.
%
\theoremstyle{plain}
\newtheorem{theorem}[subsection]{Theorem}
\newtheorem{proposition}[subsection]{Proposition}
\newtheorem{lemma}[subsection]{Lemma}

\theoremstyle{definition}
\newtheorem{definition}[subsection]{Definition}
\newtheorem{example}[subsection]{Example}
\newtheorem{exercise}[subsection]{Exercise}
\newtheorem{situation}[subsection]{Situation}

\theoremstyle{remark}
\newtheorem{remark}[subsection]{Remark}
\newtheorem{remarks}[subsection]{Remarks}

\numberwithin{equation}{subsection}

% Macros
%
\def\lim{\mathop{\mathrm{lim}}\nolimits}
\def\colim{\mathop{\mathrm{colim}}\nolimits}
\def\Spec{\mathop{\mathrm{Spec}}}
\def\Hom{\mathop{\mathrm{Hom}}\nolimits}
\def\Ext{\mathop{\mathrm{Ext}}\nolimits}
\def\SheafHom{\mathop{\mathcal{H}\!\mathit{om}}\nolimits}
\def\SheafExt{\mathop{\mathcal{E}\!\mathit{xt}}\nolimits}
\def\Sch{\mathit{Sch}}
\def\Mor{\mathop{\mathrm{Mor}}\nolimits}
\def\Ob{\mathop{\mathrm{Ob}}\nolimits}
\def\Sh{\mathop{\mathit{Sh}}\nolimits}
\def\NL{\mathop{N\!L}\nolimits}
\def\CH{\mathop{\mathrm{CH}}\nolimits}
\def\proetale{{pro\text{-}\acute{e}tale}}
\def\etale{{\acute{e}tale}}
\def\QCoh{\mathit{QCoh}}
\def\Ker{\mathop{\mathrm{Ker}}}
\def\Im{\mathop{\mathrm{Im}}}
\def\Coker{\mathop{\mathrm{Coker}}}
\def\Coim{\mathop{\mathrm{Coim}}}

% Boxtimes
%
\DeclareMathSymbol{\boxtimes}{\mathbin}{AMSa}{"02}

%
% Macros for moduli stacks/spaces
%
\def\QCohstack{\mathcal{QC}\!\mathit{oh}}
\def\Cohstack{\mathcal{C}\!\mathit{oh}}
\def\Spacesstack{\mathcal{S}\!\mathit{paces}}
\def\Quotfunctor{\mathrm{Quot}}
\def\Hilbfunctor{\mathrm{Hilb}}
\def\Curvesstack{\mathcal{C}\!\mathit{urves}}
\def\Polarizedstack{\mathcal{P}\!\mathit{olarized}}
\def\Complexesstack{\mathcal{C}\!\mathit{omplexes}}
% \Pic is the operator that assigns to X its picard group, usage \Pic(X)
% \Picardstack_{X/B} denotes the Picard stack of X over B
% \Picardfunctor_{X/B} denotes the Picard functor of X over B
\def\Pic{\mathop{\mathrm{Pic}}\nolimits}
\def\Picardstack{\mathcal{P}\!\mathit{ic}}
\def\Picardfunctor{\mathrm{Pic}}
\def\Deformationcategory{\mathcal{D}\!\mathit{ef}}
\usepackage{cleveref}
\crefname{originaltheorem}{original theorem}{original theorems}
\Crefname{originaltheorem}{Original theorem}{Original theorems}
\makeatletter
\newlabel{theorem-colimit}{{09YQ}{source}{}{sourcepointer-1}{}}
\newlabel{theorem-colimit@cref}{{[originaltheorem][0][]09YQ}{[1][0][]source}{}{}{}}
\makeatother
\crefname{originalremark}{original remark}{original remarks}
\Crefname{originalremark}{Original remark}{Original remarks}
\makeatletter
\newlabel{section-trace-method}{{03SH}{source}{}{sourcepointer-2}{}}
\newlabel{section-trace-method@cref}{{[originalremark][0][]03SH}{[1][0][]source}{}{}{}}
\makeatother
\crefname{originalproposition}{original proposition}{original propositions}
\Crefname{originalproposition}{Original proposition}{Original propositions}
\makeatletter
\newlabel{proposition-finite-higher-direct-image-zero}{{03QP}{source}{}{sourcepointer-3}{}}
\newlabel{proposition-finite-higher-direct-image-zero@cref}{{[originalproposition][0][]03QP}{[1][0][]source}{}{}{}}
\makeatother
\crefname{originallemma}{original lemma}{original lemmas}
\Crefname{originallemma}{Original lemma}{Original lemmas}
\makeatletter
\newlabel{lemma-relative-colimit}{{09Z1}{source}{}{sourcepointer-4}{}}
\newlabel{lemma-relative-colimit@cref}{{[originallemma][0][]09Z1}{[1][0][]source}{}{}{}}
\makeatother
\crefname{originalproposition}{original proposition}{original propositions}
\Crefname{originalproposition}{Original proposition}{Original propositions}
\makeatletter
\newlabel{proposition-leray}{{03QC}{source}{}{sourcepointer-5}{}}
\newlabel{proposition-leray@cref}{{[originalproposition][0][]03QC}{[1][0][]source}{}{}{}}
\makeatother
\title{Nonvanishing of degree n cohomology is inherited along a limit of finite etale covers of degree prime to l}
\author{The Stacks Project Authors}
\date{stacks-project:etale-cohomology.tex:local-snapshot-20261011; GNU FDL 1.2 or later}
\begin{document}
\maketitle
\noindent\textit{Editorial note.} All delivered mathematical text is the author's own, copied from the pinned Stacks Project chapter file etale-cohomology.tex; no statement or proof step was rewritten, reordered or removed.\par\medskip
\noindent\textit{Editorial note.} The proof is the author's complete argument for both implications, and it exhibits the transfers named in the last sentence of the statement.\par\medskip
\noindent\textit{Editorial note.} Every labelled reference occurring inside the counted spans is carried in assembly.external\_references as a non-load-bearing navigation pointer (same-chapter ones at their definition site, cross-chapter ones at the invocation inside the counted span), so no load-bearing mathematics is replaced by a stub. No bibliography citation occurs inside the counted spans.\par\medskip
\noindent\textit{Editorial note.} The checked chapter list for same-name collisions is the set of stacks chapters staged in this repository: chow, cohomology, divisors, etale-cohomology, flat, more-algebra, more-morphisms, morphisms, spaces-more-morphisms, varieties together with this card's chapter etale-cohomology.\par\medskip
\noindent\textit{Editorial note.} The preamble is the chapter's own native preamble excerpt already used for this source: the Xy-pic setup (native preamble lines 14-18) and the native theorem declarations with the mathematical macros (native preamble lines 154-217).\par\medskip
\noindent\textit{Allowed standard prerequisites (editorial).}\begin{itemize}
\item Limit of a directed system of finite etale covers of degree prime to l. Same chapter, the description of the limit X of a directed system of finite etale S-schemes of constant degree prime to l and the functors g\textasciicircum{}\{-1\}, g\_* attached to the projection g : X -> S. Supplies the morphism g and the two transfers used in the proof.
\item Behaviour of l-power torsion cohomology under finite etale covers of degree prime to l. Same chapter, the results stating that a finite etale morphism of degree prime to l induces an injection on l-power torsion cohomology, and the passage to the limit of such covers. Supplies both implications, since transfer and restriction compose to multiplication by the degree.
\end{itemize}\medskip
\ifdefined\hypertarget\hypertarget{sourcepointer-1}{}\fi
\noindent\textit{Source reference.} Original theorem 09YQ: sources/stacks-etale-cohomology/pinned-original-source.tex, source lines 7400--7400. Navigation pointer invoked at proof line 10567; the counted span contains the pointer only and the cited result is not reproduced inside it, so the pointer is registered as a non-load-bearing navigation stub in the established form for Stacks material. No bibliography citation is involved.\par\medskip
\ifdefined\hypertarget\hypertarget{sourcepointer-2}{}\fi
\noindent\textit{Source reference.} Original remark 03SH: sources/stacks-etale-cohomology/pinned-original-source.tex, source lines 10330--10330. Navigation pointer invoked at proof line 10538; the counted span contains the pointer only and the cited result is not reproduced inside it, so the pointer is registered as a non-load-bearing navigation stub in the established form for Stacks material. No bibliography citation is involved.\par\medskip
\ifdefined\hypertarget\hypertarget{sourcepointer-3}{}\fi
\noindent\textit{Source reference.} Original proposition 03QP: sources/stacks-etale-cohomology/pinned-original-source.tex, source lines 8245--8245. Navigation pointer invoked at proof line 10589; the counted span contains the pointer only and the cited result is not reproduced inside it, so the pointer is registered as a non-load-bearing navigation stub in the established form for Stacks material. No bibliography citation is involved.\par\medskip
\ifdefined\hypertarget\hypertarget{sourcepointer-4}{}\fi
\noindent\textit{Source reference.} Original lemma 09Z1: sources/stacks-etale-cohomology/pinned-original-source.tex, source lines 7535--7535. Navigation pointer invoked at proof line 10590; the counted span contains the pointer only and the cited result is not reproduced inside it, so the pointer is registered as a non-load-bearing navigation stub in the established form for Stacks material. No bibliography citation is involved.\par\medskip
\ifdefined\hypertarget\hypertarget{sourcepointer-5}{}\fi
\noindent\textit{Source reference.} Original proposition 03QC: sources/stacks-etale-cohomology/pinned-original-source.tex, source lines 8176--8176. Navigation pointer invoked at proof line 10594; the counted span contains the pointer only and the cited result is not reproduced inside it, so the pointer is registered as a non-load-bearing navigation stub in the established form for Stacks material. No bibliography citation is involved.\par\medskip
\begin{lemma}
\label{lemma-nonvanishing-inherited}
Let $\ell$ be a prime number and $n$ an integer $> 0$.
Let $S$ be a quasi-compact and quasi-separated scheme.
Let $X = \lim_{i \in I} X_i$ be the limit of a
directed system of $S$-schemes each $X_i \to S$
being finite \'etale of constant degree relatively prime to $\ell$.
The following are equivalent:
\begin{enumerate}
\item there exists an $\ell$-power torsion sheaf
$\mathcal{G}$ on $S$ such that $H_\etale^n(S, \mathcal{G}) \neq 0$ and
\item there exists an $\ell$-power torsion sheaf $\mathcal{F}$ on $X$
such that $H_\etale^n(X, \mathcal{F}) \neq 0$.
\end{enumerate}
In fact, given
$\mathcal{G}$ we can take $\mathcal{F} = g^{-1}\mathcal{F}$
and given
$\mathcal{F}$ we can take $\mathcal{G} = g_*\mathcal{F}$.
\end{lemma}

\begin{proof}
Let $g : X \to S$ and $g_i : X_i \to S$ denote the structure morphisms.
Fix an $\ell$-power torsion sheaf $\mathcal{G}$ on $S$
with $H^n_\etale(S, \mathcal{G}) \not = 0$.
The system given by $\mathcal{G}_i = g_i^{-1}\mathcal{G}$
satisfy the conditions of Theorem \ref{theorem-colimit}
with colimit sheaf given by $g^{-1}\mathcal{G}$. This tells 
us that:
$$
\colim_{i\in I} H^n_\etale(X_i, g_i^{-1}\mathcal{G}) =
H^n_\etale(X, \mathcal{G})
$$
By virtue of the $g_i$ being finite \'etale morphism of degree prime
to $\ell$ we can apply ``la m\'ethode de la trace'' and we find
the maps
$$
H^n_\etale(S, \mathcal{G}) \to H^n_\etale(X_i, g_i^{-1}\mathcal{G})
$$
are all injective (and compatible with the transition maps).
See Section \ref{section-trace-method}. Thus, the colimit is non-zero, i.e.,
$H^n(X,g^{-1}\mathcal{G}) \neq 0$, giving us the desired result with 
$\mathcal{F} = g^{-1}\mathcal{G}$.

\medskip\noindent
Conversely, suppose given an $\ell$-power torsion sheaf $\mathcal{F}$ on $X$
with $H^n_\etale(X, \mathcal{F}) \not = 0$. We note that since the $g_i$
are finite morphisms the higher direct images vanish
(Proposition \ref{proposition-finite-higher-direct-image-zero}).
Then, by applying Lemma \ref{lemma-relative-colimit}
we may also conclude the  same for $g$.
The vanishing of the higher direct images tells us that
$H^n_\etale(X, \mathcal{F}) = H^n(S, g_*\mathcal{F}) \neq 0$
by Leray (Proposition \ref{proposition-leray})
giving us what we want with $\mathcal{G} = g_*\mathcal{F}$.
\end{proof}

\end{document}
